\subsection{1D Density Pulse Advection}
\label{SEC:Testcases:DensityPulse}

\begin{figure}[h!tb]
  \begin{centering}
    \begin{subfigure}{3.0in}
      \insertimage[width=\textwidth]{test_cases/density_pulse_start}
      \caption{Initial condition.}
      \label{FIG:DensityPulseStart}
    \end{subfigure}%
    \begin{subfigure}{3.0in}
      \insertimage[width=\textwidth]{test_cases/density_pulse_final}
      \caption{\PO[3] result after advecting one pulse width.}
      \label{FIG:DensityPulseFinal}
    \end{subfigure}
    \caption{Density pulse test case.}
    \label{FIG:DensityPulse}
  \end{centering}
\end{figure}

This test case features the inviscid advection of a density pulse, and can be
used to assess the numerical dissipation of the scheme.
%
\Fref{FIG:DensityPulseStart} shows the starting density profile and the
available user input parameters for adjusting the center, width, and gap of the
profile as well as the domain length.
%
Gauss, square, triangular, and elliptical pulse shapes with an amplitude of
1.0~\unitDens{} can be superimposed onto the reference flow
conditions.
%
However, the combination pulse is recommended because it shows the effect on all
the profiles shapes in one simulation.
%
Periodic boundary conditions are applied.
%
\Fref{FIG:DensityPulseFinal} shows a sample \PO[3] result using the resolution
marker to identify troubled cells and applying a filter with
$\inp[nmlLimiter::LocalFilterOrder]{Local\_Filter\_Order}=4$ to damp
oscillations.

%\input{test_cases/density_pulse_inp}

\SeeAlso{
\inp[nmlInitialization::Initialization]{Initialization}=24,
\inp[nmlGrid::InternalGridDoF]{Internal\_Grid\_DoF},
\inp[nmlTestcasePulse]{nml\_testcase\_pulse}
}

